% ==============================================================================
% 04_mo_hinh_hoa.tex — Phần 4. Mô hình hóa phân tích
% Phụ trách: Xuân (Dev 2)
% Ngân sách trang: 6 trang
% ==============================================================================

\pttkPart{Mô hình hóa phân tích}

\todosection{Xuân (Dev 2)}{Tham chiếu \texttt{glossary\_lock.md}, \texttt{docs/ARCHITECTURE\_VI.md} và \texttt{docs/xuan/01..15}. Ngân sách mục tiêu: 6 trang.}

\subsection{Biểu đồ Use Case tổng quát}
\todofig{Hình 3. Biểu đồ Use Case tổng quát của DevFlow (UC01--UC12) — diagrams/out/hinh\_03\_usecase.pdf}

\subsection{Đặc tả Use Case trọng tâm}
% Đặc tả đầy đủ 3 UC cốt lõi: UC05 (Kéo thả & đồng bộ Board), UC06 (Git Webhook & Auto-move), UC07/UC08 (AI CI Triage / Task Breakdown).

\subsection{Biểu đồ tuần tự (Sequence Diagram)}
\todofig{Hình 4. Biểu đồ tuần tự luồng Git Webhook $\rightarrow$ Event Bus $\rightarrow$ Board \& AI — diagrams/out/hinh\_04\_sequence\_git.pdf}

\subsection{Biểu đồ trạng thái của Công việc (Task)}
\todofig{Hình 5. Máy trạng thái đối tượng Task/Card (Thủ công và Tự động từ Git) — diagrams/out/hinh\_05\_state\_task.pdf}
% Bảng chuyển trạng thái (Từ, Sự kiện, Đến, Ai/Hệ thống thực hiện, Điều kiện).

\subsection{Mô hình miền và ngữ cảnh giới hạn (DDD)}
% Bảng 5 Bounded Context của DevFlow: Auth, Board, Git & CI, AI Service, Notification.
